% Generated draft: requires mathematical and source review.
\documentclass[11pt]{article}
\usepackage[margin=25mm]{geometry}
\usepackage{fontspec}
\setmainfont{Latin Modern Roman}
\usepackage{amsmath,amssymb,mathtools,unicode-math}
\setmathfont{Latin Modern Math}
\usepackage{enumitem,needspace,xurl,hyperref}
\setlength{\parindent}{0pt}
\setlength{\parskip}{5pt}
\setlength{\emergencystretch}{3em}
\newcommand{\N}{\mathbb{N}}
\newcommand{\Z}{\mathbb{Z}}
\newcommand{\Q}{\mathbb{Q}}
\newcommand{\R}{\mathbb{R}}
\newcommand{\C}{\mathbb{C}}
\newcommand{\F}{\mathbb{F}}
\newcommand{\A}{\mathbb{A}}
\newcommand{\PP}{\mathbb{P}}
\newcommand{\E}{\mathbb{E}}
\newcommand{\eps}{\varepsilon}
\newcommand{\dd}{\,\mathrm{d}}
\begin{document}
\section*{UnsolvedMath ID: 30002199}
\textbf{Source problem code:} catalogue entry
\subsection{Definitions and mathematical statement}
Let \(T=(\C^*)^n\) be the given complex algebraic torus, let \(Y\subseteq T\) be a subvariety, and let \(X_\Sigma\) be a toric variety whose dense torus is identified with \(T\). Write
\[
\overline{Y}^{\,X_\Sigma}
\]
for the Zariski closure of \(Y\) in \(X_\Sigma\). For a point \(x\in X_\Sigma\), its \(T\)-orbit is
\[
T\cdot x=\{t\cdot x:t\in T\}.
\]
\textbf{Question.} For every such pair \((Y,X_\Sigma)\), characterize the set
\[
\mathcal{I}(Y,X_\Sigma)
=
\bigl\{O\subseteq X_\Sigma: O\text{ is a }T\text{-orbit and }O\cap\overline{Y}^{\,X_\Sigma}\ne\varnothing\bigr\}.
\]
Thus the requested characterization should determine, for each \(T\)-orbit of the given toric variety, whether it meets the closure of \(Y\). No completeness or projectivity hypothesis on \(X_\Sigma\) is supplied.

The supplied curation describes this as a question extracted from a tropicalization passage, but states that the passage does not preserve a distinct conjecture or theorem to verify. Accordingly, this draft records the orbit-intersection question without asserting a proposed criterion or an independently verified open-problem status.
\subsection{Short English statement}
Given a subvariety of a complex algebraic torus and a toric variety containing that torus as its dense orbit, determine exactly which torus orbits meet the subvariety's Zariski closure. The supplied record labels the question open, while its dated assessment says that no distinct conjecture or theorem was preserved.
\subsection{Sources}
\begin{itemize}
\item Supplied local registry record: permanent UnsolvedMath ID \(30002199\), title \emph{Toric Orbits Meeting Subvariety Closures}, source problem code \texttt{catalogue entry}, collection \emph{Oberwolfach Reports Open Problems}. No individual proposer is supplied; the proposed year is recorded as 2012.
\item Supplied report citation: \emph{Mini-Workshop: Cohomology Rings and Fundamental Groups of Hyperplane Arrangements, Wonderful Compactifications, and Real Toric Varieties} (2013), \emph{Oberwolfach Reports}, pp.~2939--2983, DOI \texttt{10.4171/owr/2012/49}. Supplied source URL and literature-source link, labeled ``original OWR report'': \url{https://doi.org/10.4171/owr/2012/49}. The background describes Oberwolfach Reports as documenting mathematical workshops and their open-problem discussions.
\item Supplied literature assessment, checked \texttt{2026-08-21}: status ``open''; the extracted tropicalization passage supplies background and a question about toric-orbit intersections, without preserving a distinct conjecture or theorem to verify. This is supplied dated curation, not independently verified here. The supplied background requests independent verification of time-sensitive status and cited claims.
\item The registry labels the statement \texttt{corrected\_verified} and reports that its mathematical content was checked against the official Oberwolfach report, with notation, context, and question form normalized without changing the source-level claim. This verification is reported by the supplied record and has not been independently reproduced here.
\end{itemize}
\end{document}
